\newpage
\section{INTRODUCCIÓN}

En los procesos de manufactura moderna, el mecanizado por arranque de viruta mediante tornos de Control Numérico Computarizado (CNC) representa una de las operaciones más críticas y demandadas para la fabricación de componentes de alta precisión. Dentro de este proceso, la herramienta de corte es el elemento primario que interactúa directamente con el material de trabajo, estando sometida continuamente a severos esfuerzos mecánicos, elevadas temperaturas y fricción dinámica. Como consecuencia natural de esta interacción termomecánica, los insertos experimentan una degradación progresiva en la geometría de su filo.

La evolución del desgaste de las herramientas incide de forma determinante en la estabilidad del proceso de corte, modificando las fuerzas de mecanizado y alterando directamente la calidad del acabado superficial y la precisión dimensional de las piezas manufacturadas. Tradicionalmente, la gestión del reemplazo de insertos en la industria se ha sustentado en inspecciones visuales periódicas o en tablas estandarizadas de vida útil. Sin embargo, estas estrategias empíricas resultan inadecuadas para anticipar fallas dinámicas o desgastes anómalos generados por variaciones en las propiedades del material o fluctuaciones en las condiciones de corte.

Frente a estas limitaciones operativas, el desarrollo de sistemas de Monitoreo de Condición de Herramientas (\textit{Tool Condition Monitoring}, TCM) se ha consolidado como una disciplina fundamental dentro de la manufactura inteligente. El monitoreo indirecto y continuo permite inferir en tiempo real el estado de la herramienta a través de señales físicas del proceso, como las variaciones en la corriente del motor del husillo y las oscilaciones mecánicas de alta frecuencia, facilitando una transición hacia estrategias avanzadas de mantenimiento predictivo.

En este contexto tecnológico, el presente proyecto de grado tiene como propósito diseñar e implementar un sistema periférico no invasivo de monitoreo de condición para tornos CNC, integrando adquisición de señales de corriente y vibración, procesamiento inteligente en el borde (\textit{Edge AI}) y conectividad industrial. La investigación y validación experimental del sistema se desarrollan específicamente en el entorno operativo de la empresa metalmecánica Maestranza San Carlos (MAINSA), adaptando la solución a las exigencias reales de su línea de producción.
